\documentclass[../../main.tex]{subfiles}
\begin{document}

{\bf [解]}
$h(x)=f(x)-g(x)$ とおく．$h(x)$ は4次のモニック多項式であり，$h(1)=h(-1)=0$ より，因数定理から，$\alpha,\beta\in\mathbb{R}$ として
\begin{align}
h(x)=(x+1)(x-1)(x^2+\alpha x+\beta) \label{eq:1}
\end{align}
とかける．したがって，
\begin{align}
\{h(x)\}^2=(x^2-1)^2(x^4+2\alpha x^3+(\alpha^2+2\beta)x^2+2\alpha\beta x+\beta^2)
\end{align}
だから，$A=\displaystyle\int_{-1}^1\{h(x)\}^2dx$ として偶関数と奇関数の性質を用いると
\begin{align}
  A
   &=2\int_0^1(x^2-1)^2(x^4+(\alpha^2+2\beta)x^2+\beta^2)dx \\
\therefore\ 
 \dfrac{1}{2}A
  &=\left(\alpha^2+2\beta\right)\int_0^1(x^2-1)^2x^2dx+\beta^2\int_0^1(x^2-1)^2dx+\int_0^1x^4(x^2-1)^2dx \\
  &=\dfrac{8}{105}\alpha^2+\dfrac{8}{15}\beta^2+\dfrac{16}{105}\beta+\int_0^1x^4(x^2-1)^2dx \\
  &=\dfrac{8}{105}\alpha^2+\dfrac{8}{15}\left\{\left(\beta+\dfrac{1}{7}\right)^2-\dfrac{1}{49}\right\}+\int_0^1x^4(x^2-1)^2dx \label{eq:2}
\end{align}
である．よって$A$の最小値を与える $(\alpha,\beta)$ は，
\begin{align}
 (\alpha,\beta)=\left(0,-\dfrac{1}{7}\right) \label{eq:3}
\end{align}
である．この時\cref{eq:1}から
\begin{align}
  h(x)=x^4+\alpha x^3+(\beta-1)x^2-\alpha x-\beta = x^4-\dfrac{8}{7}x^2+\dfrac{1}{7}
\end{align}
だから，$(a,b,c,d)=\left(0,\dfrac{8}{7},-1,-\dfrac{1}{7}\right)$ とおけば良い．

{\bf [解説]}
今回出てきた積分はベータ関数
\begin{align}
 B(x,y) = \int_0^{1}t^{x-1}(1-t)^{y-1} dt
\end{align}
の形をしている．これはガンマ関数と
\begin{align}
 B(x,y) = \dfrac{\Gamma(x)\Gamma(y)}{\Gamma(x+y)}
\end{align}
の関係にある．ガンマ関数は引数が自然数の時は階乗だから$n,m\in\mathbb{N}$に対して
\begin{align}
 B(n,m) = \dfrac{(n-1)!(m-1)!}{(n+m-1)!}
\end{align}
である．有名な$1/6$公式はこの$x=y=2$の時で
\begin{align}
 B(2,2) = \dfrac{1}{3!} = \dfrac{1}{6}
\end{align}
である．

\newpage

\end{document}
